\subsubsection{Training EBMs}\label{sec:ebm_training_2d}

Next we consider the problem of training \glspl*{ebm} on a collection of synthetic two-dimensional datasets given in Appendix \ref{app:training_ebm_2d}. For each dataset, we parametrise the model energy with a neural network, which we train by minimising the negative log-likelihood of the data. Assume the model is normalisable of the form,
\begin{equation*}
    p_\theta(x)=Z_\theta^{-1}e^{-E(\theta,x)},\qquad Z_\theta=\int_{\mathbb{R}^{d_x}}e^{-E(\theta,x)}\md x,
\end{equation*}
where $E(\theta,x)$ is the energy function parametrised by $\theta\in\mathbb{R}^{d_\theta}$ and $Z_\theta$ is the normalisation constant. Given a dataset of $M$ observations denoted $\{y^m\}_{m=1}^M$, the gradient of the negative log-likelihood is given,
\begin{equation*}
    \nabla_\theta\ell(\theta) = \mathbb{E}_{\pi_\theta}\left[-\nabla_\theta E(\theta,X)+\frac{1}{M}\sum_{m=1}^M\nabla_\theta E(\theta, y^m)\right].
\end{equation*}
% Replacing the fixed data averaging with an expectation over an underlying data distribution $p_\text{data}$ makes this equivalent to minimising the (forward) \gls*{kl} divergence between the model distribution, $\pi_\theta$, and $p_{\text{data}}$ with
% \begin{equation}
%     \ell(\theta)=\mathrm{KL}(p_{\text{data}}\Vert \pi_\theta).
% \end{equation}
Here, we can see the \gls*{pcd} algorithm is a special case of the \gls*{impdiff} framework. We tune each method separately to showcase a fair comparison using separately selected hyperparameters. We evaluate performance with respect to the number of outer iterations, which is equivalent to fixing the number of gradient back-propagations using our discretisation, reporting the Fisher Divergence in Figure \ref{fig:training_comparisons}. Even with separate tuning, there is a clear acceleration of \gls*{mmfpo} in the logarithmic scale. We also report the learnt probability densities in Figure \ref{fig:models} and see both methods produce models of visually similar qualities.